

\section{Introduction}
\label{sec:intro}
\emph{Fixed points} abound in mathematics and computer science.
In logic we may enrich languages by `positive' fixed points to perform (co)inductive reasoning, while in
programming languages positive fixed points in type systems are used to represent (co)datatypes and carry out (co)recursion.
In both settings the underlying systems may be construed as fragments of their second-order counterparts.
% However the intuitionism underlying type systems allows for alternative semantic interpretations too, not least as bona fide fixed points themselves.

In this work we investigate the computational expressivity of type systems with least and greatest (positive) fixed points.
We pay particular attention to \emph{circular proof systems}, where typing derivations are possibly non-well-founded (but regular), equipped with an $\omega$-regular `correctness criterion' at the level of infinite branches.
Such systems have their origins in modal fixed point logics, notably the seminal work of Niwi\'nski and Walukiewicz \cite{NiwWal96:games-for-mu-calc}.
Viewed as type systems under the `Curry-Howard' correspondence, circular proofs have received significant attention in recent years, notably based in systems of \emph{linear logic} \cite{BaeldeDS16Infinitaryprooftheory,EhrJaf21:cat-models-muLL,EhrJafSau21:totality-CmuMALL,BDKS22:bouncing-threads,DeSau19:infinets1,DePelSau:infinets2,DeJafSau22:phase-semantics} after foundational work on related finitary systems in \cite{BaeldeMiller07,Baelde12}.
In these settings circular proofs are known to be (at least) as expressive as their finitary counterparts, but classifying the exact expressivity of both systems has remained an open problem.
This motivates the main question of the present work:

\smallskip

\begin{question}
\label{question:main}
What functions do (circular) proof systems with fixed points represent?
\end{question}


\medskip


Circular type systems with fixed points were arguably pre-empted by foundational work of Clairambault \cite{Clairambault09}, who introduced an extension $\muLJ$ of Gentzen's sequent calculus $\LJ$ for intuitionistic propositional logic by least and greatest positive fixed points.
$\muLJ$ forms the starting point of our work and, using standard methods, admits an extension into a circular calculus, here called $\cmuLJ$, whose computational content we also investigate.

% \medskip

In parallel lines of research, fixed points have historically received considerable attention within mathematical logic.
% , not least in descriptive set theory.
The ordinal analysis of extensions of Peano Arithmetic ($\PA$) by inductive definitions has played a crucial role in giving proof theoretic treatments to (impredicative) second-order theories (see, e.g., \cite{sep-proof-theory}).
More recently, inspired by Lubarsky's work on `$\mu$-definable sets' \cite{Lubarsky93}, M\"ollerfeld has notably classified the proof theoretic strength of extensions of $\PA$ by general inductive definitions in \cite{Moellerfeld02:phd-thesis}.

In this work we somewhat bridge these two traditions, in computational logic and in mathematical logic, in order to answer our main question.
In particular we apply proof theoretic and metamathematical techniques to show that both $\muLJ$ and $\cmuLJ$ represent precisely the functions provably recursive in the subsystem $\PSCA$ of second-order arithmetic.
This theory is far beyond the `big five' of reverse mathematics, and is among the strongest theories for which we have an effective ordinal analysis (see \cite{Rathjen95:recent-advances}). 
The best known lower bound for $\muLJ$ before was G\"odel's $\T$ (see, e.g., \cite{Clairambault09}), which has the same proof theoretic strength as $\PA$. 
The best known upper bound was Girard-Reynold's $\F$, thanks to its impredicative encodings of fixed points, which has the same proof theoretic strength as second-order arithmetic $\SOPA$.


% Thus this work resolves our main Question~\ref{question:main}.

\subsection{Outline and contribution}
The structure of our overall argument is visualised in Figure~\ref{fig:grand-tour}, outlining a cycle of inclusions of `representable functions'.
Here the upper row consists of theories of arithmetic, where the representable functions of an arithmetic theory $T$ are just its provably total recursive functions; i.e.\ those functions $ f : \Nat \to \cdots \to \Nat$ with graph computed by some $\Sigma^0_1$ formula $\phi_f(\vec x, y)$ such that $T \proves \forall \vec x \exists y \phi_f (\vec x,y) $.
The lower row consists of type systems whose representable functions are just those admitting a typing derivation with conclusion $\nat \to \cdots \to \nat$ computing the function under its operational semantics (as in, e.g., Definition~\ref{defn:numerals-representability}). 

(1) is a standard embedding of finitary proofs into circular proofs (Proposition~\ref{prop:cmulj-simulates-mulj}).
(2) reduces $\cmuLJ$ to its `negative fragment', in particular free of greatest fixed points ($\nu$), via a double negation translation (Proposition~\ref{prop:cmulj-into-cmuljminus}). 


(3) is one of our main contributions: we build a higher-order computability model $\HR$ that interprets $\cmuLJneg$ (Theorem~\ref{thm:prog-implies-hr}), and moreover formalise this construction itself within $\PSCA$ to obtain our upper bound (Theorem~\ref{thm:cmuLJneg-to-psca}).
The domain of this model a priori is an (untyped) term extension of $\cmuLJneg$.
It is important for logical complexity that we interpret fixed points semantically as bona fide fixed points, rather than via encoding into a second-order system.
Along the way we must also establish some novel reverse mathematics of the Knaster-Tarski fixed point theorem (Theorem~\ref{thm:muphi-is-phiwo}).

(4) is an intricate and nontrivial result established by M\"ollerfeld in \cite{Moellerfeld02:phd-thesis}, which we use as a `black box'.
(5) is again a double negation translation, morally a specialisation of the $\Pin 0 2 $-conservativity of full second-order arithmetic $\SOPA$ over its intuitionistic counterpart $\SOHA$, composed with a relativisation of quantifiers to $\Nat$ (\Cref{thm:muPA-pi02-cons-muHA,prop:muHA-to-muHAneg}). 

(6) is our second main contribution: we provide a realisability interpretation from $\muHAneg$ into $\muLJ$ (Theorem~\ref{thm:realisability-soundness}), morally by considerable specialisation of the analogous interpretation from $\SOHA$ into Girard-Reynolds' system $\F$.
Our domain of realisers is a (typed) term extension of $\muLJneg$ (the negative fragment of $\muLJ$), which is itself interpretable within $\muLJ$ (Proposition~\ref{prop:cmulj-into-cmuljminus}).

% Finally, the application of our results to characterise the representable functions of (circular) $\muMALL$ is given in Section~\ref{sec:further} (Theorem~\ref{thm:(c)mulj-(c)mumall}).



\subsection{Related work}
Fixed points have been studied extensively in type systems for programming languages.
In particular foundational work by Mendler in the late '80s \cite{Mendler87:recursive-types,Mendler91:inductive-types} already cast inductive type systems as fragments of second-order ones such as Girard-Reynolds' $\F$ \cite{girard1972systemF,Reynolds74systemF}.
Aside from works we have already mentioned, (a variant of) (5) has already been obtained by Tupailo in \cite{Tupailo04doubleneg}.
Berger and Tsuiki have also obtained a similar result to (6) in a related setting \cite{BerTsu21:ifp}, for \emph{strictly} positive fixed points, where bound variables may never occur under the left of an arrow. 
Their interpretation of fixed points is more akin to that in our type structure $\HR$ than our realisability model.

Finally the structure of our argument, cf.~Figure~\ref{fig:grand-tour}, is inspired by recent works in cyclic proof theory, notably \cite{Simpson17:cyc-arith,Das20:ca-log-comp} for (cyclic) (fragments of) $\PA$ and \cite{Das21:CT-preprint,Das21:CT-fscd,KuperbergPP21systemT} for (circular) (fragments of) G\"odel's system $\T$.


\begin{figure}
    % \centering
    \begin{tikzpicture}
% \node (clo-muLJneg) {$\closure\muLJneg$};
% \node [right = of clo-muLJneg] (muLJ) {$\muLJ^{(-)}$};
% \node [right = of muLJ] (cmuLJ) {$\cmuLJ$};
% \node [right = of cmuLJ] (cmuLJneg) {$\cmuLJneg$};
% \node [right = of cmuLJneg] (clo-cmuLJneg) {$\closure \cmuLJneg$};
% \node [above = of clo-muLJneg] (muHA) {$\muHA$};
% \node [above = of clo-cmuLJneg] (PSCA) {$\PSCA$};
% \node [above = of cmuLJ] (muPA) {$\muPA$};
% \draw[->] (clo-muLJneg) to (muLJ);
% \draw[->] (muLJ) to (cmuLJ);
% \draw[->] (cmuLJ) to (cmuLJneg);
% \draw[->] (cmuLJneg) to (clo-cmuLJneg);
% \draw[->] (clo-cmuLJneg) to (PSCA);
% \draw[->] (PSCA) to (muPA);
% \draw[->] (muPA) to (muHA);
% \draw[->] (muHA) to (clo-muLJneg);
% \node [left = of clo-muLJneg] (models) {Computational models:};
% \node [above = of models] (theories) {Arithmetic theories:};
\node (muLJ) {$\muLJ^{(-)}$};
\node [right = of muLJ] (cmuLJ) {$\cmuLJ^{\phantom{-}}$};
\node [right = of cmuLJ] (cmuLJneg) {$\cmuLJneg$};
\node [above = of muLJ] (muHA) {$\muHAneg$};
\node [above = of cmuLJneg] (PSCA) {$\PSCA$};
\node [above = of cmuLJ] (muPA) {$\muPA$};
\draw[->] (muLJ) to node[below]{\small (1)} (cmuLJ);
\draw[->] (cmuLJ) to node[below]{\small (2)} (cmuLJneg);
\draw[->] (cmuLJneg) to node[right]{\small (3)} (PSCA);
\draw[->] (PSCA) to node[above]{\small (4)} (muPA);
\draw[->] (muPA) to node[above]{\small (5)} (muHA);
\draw[->] (muHA) to node[left]{\small (6)} (muLJ);
\node [left = .5cm of muLJ, align=center] (models) {\footnotesize Type \\ \footnotesize systems:};
\node [left = .5cm of muHA, align=center] (theories) {\footnotesize Arithmetic \\ \footnotesize theories:};
    \end{tikzpicture}
    \caption{Summary of the main `grand tour' of this work. All arrows indicate inclusions of representable functions.
    }
    \label{fig:grand-tour}
\end{figure}

\subsection{Comparison to preliminary version}
This paper is an expansion of the preliminary conference version \cite{CD23:muLJ-lics}.
In this version we additionally include full proofs of all our results, as well as further examples and narrative. 

We have reformulated our realisability argument in \Cref{sec:realisability} into a form of \emph{abstract} realisability, inspired by the approach of \cite{BerTsu21:ifp}. 
This factors the approach of the preliminary version by a more careful relativisation of quantifiers to deal with an inconvenient type mismatch when realising inductive predicates.

Finally, in the preliminary version we also showed equivalence of $\muLJ$ and $\cmuLJ$ with their counterparts in linear logic from (see, e.g.,~\cite{Baelde12,BaeldeDS16Infinitaryprooftheory}) via appropriate proof interpretations.
These results will be expanded upon in a separate self-contained paper.

\subsection{Notation}
% Due to space constraints, most of our proofs and technical development is devolved to appendices. 
% 
Throughout this work we employ standard rewriting theoretic notation. 
Namely for a relation $\rightsquigarrow_a$, we denote by $\rightsquigarrow^*_a$ the reflexive and transitive closure of $\rightsquigarrow_a$, and by $=_{a}$ the relexive symmetric transitive closure of $\rightsquigarrow_a$.

We shall make use of \emph{(first-order) variables}, written $x,y$ etc., and \emph{(second-order) variables}, written $X,Y$ etc.\ throughout.
We shall use these both in the setting of type systems and arithmetic theories, as a convenient abuse of notation.

